% =====================================================================
% PROPOSED (assistant draft) -- NOT Pete's paper.
%
% Written by the paper drafting and citations module as a reference draft of
% a complete paper, for comparison with Pete's own working draft. It is never
% merged into his text except where he chooses.
%
% Format follows paper/JON-RULES.md (verified 12 Oct 2026): single column,
% 12 pt Times New Roman, 2.54 cm margins, UK spelling, headings in capitals,
% numbered paragraphs 1., 1.1., equations numbered from (1), table captions
% above, figure captions below, Harvard references with no in-text page
% numbers except for locators that do real work.
%
% Section budget and figure list: see README.md in this directory.
% Factual source: paper/REFERENCE-DRAFT.md only. Every number carries the
% label its source note attached. Unverified numbers appear as [VERIFY].
% =====================================================================
\documentclass[12pt,a4paper]{article}
\usepackage[utf8]{inputenc}
\usepackage[T1]{fontenc}
\usepackage{mathptmx}                     % Times for text and maths
\usepackage[a4paper,margin=2.54cm]{geometry}
\usepackage{amsmath,amssymb,graphicx,booktabs,array}
\usepackage[hidelinks]{hyperref}
\usepackage{fancyhdr}
\usepackage{setspace}
\singlespacing
\setlength{\parindent}{0pt}
\setlength{\parskip}{6pt}
\pagestyle{fancy}\fancyhf{}
\fancyhead[L]{\small PROPOSED (assistant draft) --- not for submission}
\fancyhead[R]{\small\thepage}
\renewcommand{\headrulewidth}{0.4pt}

% --- notation shorthands -------------------------------------------------
\newcommand{\Ev}{\mathcal{Y}}
\newcommand{\state}{\mathbf{x}}
\newcommand{\impact}{\boldsymbol{\psi}}
\newcommand{\traj}{\mathbf{e}}
\newcommand{\nuis}{\boldsymbol{\theta}}

\begin{document}

% =====================================================================
% TITLE BLOCK  (JON: short title <= 40 characters)
% =====================================================================
\begin{center}
{\large\bfseries An Integrated Bayesian Reconstruction of the Flight Path,\\
Impact and Seabed Location of MH370}\\[4pt]
{\small Short title: Integrated Bayesian search for MH370 \quad(36 characters)}\\[10pt]
{P.\ O.\ Large\textsuperscript{1}}\\[4pt]
{\small\textsuperscript{1}\,[AFFILIATION]\quad(corresponding author: [EMAIL]; ORCID required at submission)}
\end{center}

\vspace{6pt}
\hrule
\vspace{6pt}

% =====================================================================
% ABSTRACT  (JON: self-contained, single paragraph, about 150 words)
% =====================================================================
\noindent\textbf{Abstract.}
The 2014 loss of Malaysia Airlines flight MH370 has been reconstructed by Bayesian
estimation from satellite communication metadata, but published reconstructions stop at the
final transmission and treat fuel, the end of flight and the subsequent fate of the aircraft
outside the estimator. This paper reports an integrated reconstruction that carries one
probability distribution from the last radar fix to the seabed. A particle filter on the
published measurement models is extended with an inside-the-likelihood fuel and endurance
model, primary radar, a waypoint route family and a modified sampler; its posterior is handed
to a terminal stage that simulates descent to impact; and the impact distribution is then
re-weighted by ocean drift of recovered debris, unsuccessful seabed search, satellite imagery
and hydroacoustic prediction. Unsuccessful search is the largest single factor. No result is
yet converged, and the paper reports the diagnostics that show it.

\vspace{4pt}
\noindent\textbf{Keywords.} (Selected at submission from the journal's list; maximum four.
Candidates: Bayesian estimation; satellite navigation; search and rescue; aircraft accident.)

% =====================================================================
\section*{1.\quad INTRODUCTION}
\addcontentsline{toc}{section}{Introduction}
% =====================================================================

Malaysia Airlines flight MH370 left Kuala Lumpur for Beijing on 7 March 2014 and did not
arrive. The aircraft, a Boeing 777-200ER registered 9M-MRO, ceased secondary surveillance
reporting less than an hour after departure, was tracked by primary radar on a reversal back
across the Malay peninsula, and then flew for a further six hours while exchanging periodic
handshakes with an Inmarsat geostationary satellite. Those handshakes carry two
quantities --- a burst timing offset, which constrains range to the satellite, and a burst
frequency offset, which constrains the Doppler shift and therefore velocity. They are the only
observations of the aircraft after 18:22~UTC. Despite three underwater searches, the principal
wreckage has not been found.

The reconstruction problem has a specific and unusual shape. The observations are few, weak
and indirect; the unknown is a six-hour trajectory; and the quantity of interest --- where the
aircraft struck the water --- lies beyond the last observation. Bayesian estimation is the
natural response, and Davey et al.\ (2016) is the authoritative treatment: a particle filter
over aircraft state, driven by the handshake data, with a prior formed from the primary radar
track and a manoeuvre model of random turns and speed changes. Its posterior at the last
transmission defined the underwater search areas.

That reconstruction stops where the data stop, and three consequences follow. First, fuel
enters only as an after-the-fact censor on trajectories rather than as a constraint inside the
likelihood, so trajectories that the aircraft could not have flown on its fuel load are
weighted and then discarded instead of being down-weighted while the filter runs. Second, the
two transmissions at 00:19~UTC --- which occurred after a power interruption and are the only
evidence of the aircraft's final minutes --- are used unevenly, and the descent they may imply
is not propagated. Third, everything that happened after impact is treated separately:
drift of the recovered debris, the unsuccessful searches, the satellite imagery and the
hydroacoustic record are each analysed against a reconstruction they cannot inform.

This paper reports an integrated reconstruction that removes those three separations. One
probability distribution is carried from the last radar fix at 18:01~UTC through to a
distribution over resting places on the seabed. The flight-path estimator is extended rather
than replaced: the published measurement models are retained and cited, and the paper
describes only what was changed or added. Fuel and endurance enter the likelihood; primary
radar enters as a likelihood term rather than only as a prior; a waypoint route family is
added to the proposal; the sampler is modified where the posterior is multi-modal; the
treatment of the 00:19 transmissions is enumerated rather than chosen; the end of flight is
simulated to impact under competing descent families; and the impact distribution is then
re-weighted, on one common set of impact samples, by every subsequent line of evidence.

Three results of that integration are reported. Unsuccessful seabed search is the largest
single factor acting on the impact distribution, and it moves the distribution north and makes
it larger rather than smaller. Ocean drift of the recovered debris widens the distribution
instead of narrowing it, and does not discriminate between the available treatments of the
00:19 data. And the hydroacoustic record, which is often cited as either confirming or
excluding particular impact locations, is at present predictive only: its log-likelihood ratios
lie within noise, for reasons of calibration rather than of physics, and the paper quantifies
what would change if that calibration existed.

A fourth result is methodological and is reported as plainly as the others. None of the
composed distributions is converged. The diagnostics that establish this --- split-half
agreement between balanced partitions of independent replicates, and effective sample size per
evidence factor --- are reported for every product, and the paper states which conclusions
survive and which do not. Published work on this accident states no quantitative convergence
criterion at all; this paper proposes one and then fails it.

% =====================================================================
\section*{2.\quad THE PUBLISHED RECONSTRUCTION, AND WHAT THIS PAPER CHANGES}
% =====================================================================

The measurement models are taken unchanged from Ashton et al.\ (2015), which is the published
account of the satellite communication analysis and the source of the satellite ephemeris, the
burst timing and frequency offset definitions, and their error models. The estimator is taken
from Davey et al.\ (2016), whose chapter structure allows each element to be cited precisely:
the measurement model in chapter~5, the cruise dynamics in chapter~6, the manoeuvre dynamics in
chapter~7, and the particle filter implementation in chapter~8. Rather than restate that
machinery, this paper follows the practice of citing the originating work by its own structure
and develops only the differences, in the manner of Ng et al.\ (2024).

Four properties of the published estimator matter for what follows. It models cruise as an
Ornstein--Uhlenbeck process about a commanded Mach number and track, with manoeuvres arriving
as a Poisson process whose rate is set by a manoeuvre time constant drawn from a broad prior.
It assumes infinite fuel inside the filter, applying fuel as a censor afterwards
(Davey et al., 2016, assumption 4). It carries no vertical rate in the cruise state, so the
frequency offset residual attributable to vertical motion is not represented. And it uses both
burst timing offsets at 00:19 but neither frequency offset, applying a log-on correction to the
anomalous second timing offset.

The factual record of the flight and of the recovered debris is the Malaysian safety
investigation report (Malaysian ICAO Annex 13 Safety Investigation Team for MH370, 2018), and
the searched areas and their coverage are reported by the Australian Transport Safety Bureau
(2017).

The reproduction of that work is exact and is the baseline of this paper. A configuration
holding the published priors and models reproduces the published log evidence to the last
stored digit, and is regression-tested so that it continues to do so as extensions are added.
Every result that depends on an extension names the extension. The two places where the
published text contradicts itself are documented in the project record and resolved
explicitly rather than silently.

Table~1 classifies every element of the integrated estimator against the published work and
states its build status. The classification is used consistently in what follows: an
\emph{update} is the same model element supplied with better data; a \emph{modification} is the
same element with a changed assumption or method; an \emph{extension} is an element the
published work does not contain.

% =====================================================================
\section*{3.\quad THE INTEGRATED ESTIMATOR}
% =====================================================================

\subsection*{3.1.\quad \textit{Unknowns and evidence}}

Let $K$ denote 00:11~UTC, the last handshake before the power interruption, and let $I$ denote
the time of water impact. The evidence is partitioned by the stage of the problem at which it
acts:
%
\begin{equation}
\Ev = \bigl\{\,\Ev^{-},\ \Ev^{19},\ \Ev^{D},\ \Ev^{S},\ \Ev^{P},\ \Ev^{A}\,\bigr\},
\end{equation}
%
where $\Ev^{-}$ is everything observed up to and including $K$ --- the primary radar track to
18:22, and the burst timing and frequency offsets of the handshakes --- $\Ev^{19}$ the two
transmissions at 00:19 together with the log-on timing that accompanies them, $\Ev^{D}$ the
identity, condition and discovery place and time of the recovered debris together with the
coastlines on which nothing was recovered, $\Ev^{S}$ the coverage and detection performance of
the unsuccessful underwater searches, $\Ev^{P}$ satellite imagery of candidate floating
objects, and $\Ev^{A}$ the hydroacoustic record of the International Monitoring System.

The unknowns are of four kinds. The pre-impact aircraft state $\state_t$ carries position,
altitude, three-dimensional velocity, mass and fuel remaining, commanded Mach number and track,
and the control states needed to propagate them. The terminal trajectory
$\traj_{K:I}$ is the complete path from $K$ to impact, including engine, electrical and
fuel-exhaustion transitions. The \emph{impact complex} $\impact$ collects what the subsequent
evidence actually depends on: first-contact time and position, mass, three-dimensional velocity,
attitude, total and vertical kinetic energy, the duration over which that energy is deposited,
and a debris class. And $\nuis$ collects continuous nuisance parameters --- frequency offset
bias and its drift, oscillator behaviour, fuel-performance corrections, winds, drift
coefficients, acoustic source and propagation parameters.

Discrete model choices are kept explicit and are not averaged away. Write
$H = (s, h, d, c, m)$ for the combination of a core stratum $s$ (which route or flight-mode
family the trajectory belongs to), an end-of-flight family $h$, a treatment $d$ of the 00:19
data, a cause $c$ of the 00:19 power cycle, and an ocean model $m$. Each is a declared
alternative with a prior, and alternatives shared between modules are aligned by name so that
they are marginalised jointly rather than independently.

\subsection*{3.2.\quad \textit{The joint posterior}}

Collecting the conditional structure gives the posterior from which every result in this paper
is a marginal:
%
\begin{equation}
\begin{split}
p\bigl(\state_{0:K},\traj_{K:I},\impact,H,\nuis \mid \Ev\bigr) \ \propto\ &
 p(H)\,p(\nuis \mid H)\,p(\state_0 \mid \nuis, H)
 \prod_{k=1}^{K} p\bigl(\state_k \mid \state_{k-1}, \nuis, H\bigr)\\[2pt]
&\times\ L^{-}\bigl(\Ev^{-} \mid \state_{0:K}, \nuis\bigr)\\[2pt]
&\times\ p\bigl(\traj_{K:I} \mid \state_K, \nuis, H\bigr)\,
          \delta\bigl(\impact - \Psi(\traj_{K:I})\bigr)\,
          L_{19}\bigl(\Ev^{19} \mid \traj_{K:I}, \nuis, H\bigr)\\[2pt]
&\times\ L_{D}\bigl(\Ev^{D} \mid \impact, \nuis, H\bigr)\,
          L_{S}\bigl(\Ev^{S} \mid \impact\bigr)\,
          L_{P}\bigl(\Ev^{P} \mid \impact, \nuis, H\bigr)\,
          L_{A}\bigl(\Ev^{A} \mid \impact, \nuis\bigr),
\end{split}
\end{equation}
%
in which $\Psi(\cdot)$ is the deterministic map from a terminal trajectory to its impact
complex, and the delta function records that $\impact$ is a function of $\traj_{K:I}$ rather
than an independent unknown.

Equation~(2) is a statement of the model, not a description of the algorithm. Sampling it
jointly would be wasteful, because the factorisation has a property worth exploiting: the
post-impact evidence $\Ev^{D}$, $\Ev^{S}$, $\Ev^{P}$ and $\Ev^{A}$ depends on the trajectory
\emph{only} through $\impact$. Conditional on the impact complex, those four likelihoods are
independent of everything upstream. The estimator is therefore organised as four stages, and
the four factors on the last line of~(2) are evaluated once each on a single shared set of
impact samples.

\subsection*{3.3.\quad \textit{Stage 1: the flight-path filter}}

Stage~1 targets the posterior over trajectory and nuisance parameters given the evidence to
$K$, within a stratum:
%
\begin{equation}
p\bigl(\state_{0:K}, \nuis \mid \Ev^{-}, s\bigr)
 \ \propto\ p(\nuis \mid s)\, p(\state_0 \mid \nuis, s)
   \prod_{k=1}^{K} p\bigl(\state_k \mid \state_{k-1}, \nuis, s\bigr)\,
   \prod_{k} L\bigl(y_k \mid \state_k, \nuis\bigr),
\end{equation}
%
approximated by a weighted particle set $\{\state_{0:K}^{(i)}, \nuis^{(i)}, w^{(i)}\}_{i=1}^{N}$
with the dynamics of Davey et al.\ (2016, chapters~6 and~7) and the measurement models of
Ashton et al.\ (2015). The stratum evidence follows as the normalising constant,
%
\begin{equation}
Z_s \;=\; \sum_{i=1}^{N} w^{(i)},
\end{equation}
%
and is the quantity by which strata are later mixed and by which discrete alternatives are
compared.

Three things differ from the published filter. Fuel and endurance are evaluated inside the
likelihood, so a trajectory whose fuel burn is inconsistent with the known load is
down-weighted as the filter advances rather than discarded afterwards. The primary radar track
contributes a likelihood term in addition to shaping the prior. And the cruise state carries a
vertical rate, so that the frequency offset residual attributable to vertical motion is
represented rather than set to zero. The last of these was expected to matter and does not: its
effect on the posterior is an order of magnitude smaller than the noise against which it was
being judged, and that is reported in section~5 as a finding rather than suppressed as a
disappointment.

\subsection*{3.4.\quad \textit{The hand-off}}

Stage~1 stops at $K$ and its posterior is resampled to form the initial distribution of the
terminal stage. Writing $M$ for the number of hand-off rows,
%
\begin{equation}
q\bigl(\state_K \mid \Ev^{-}, s\bigr) \;=\; \frac{1}{M}\sum_{j=1}^{M}
  \delta\bigl(\state_K - \state_K^{(j)}\bigr),
  \qquad \state_K^{(j)} \sim \bigl\{\state_K^{(i)}, w^{(i)}\bigr\},
\end{equation}
%
each row carrying the full state and the nuisance parameters that generated it, so that the
terminal stage inherits the correlations the filter established rather than a summary of them.

The hand-off makes a distinction that must be stated because it is easy to lose. A run that
writes a hand-off must stop \emph{before} the transmissions the terminal stage will handle; its
posterior is therefore at 00:11, not at 00:19. A run that scores the 00:19 transmissions in the
filter cannot write a hand-off. The two are different estimates of different things and are not
interchangeable: the first is what should be quoted for the position at the last transmission,
the second is what feeds the terminal stage.

\subsection*{3.5.\quad \textit{Stage 2: end of flight, and the 00:19 transmissions}}

From each hand-off row the terminal stage propagates powered flight to fuel exhaustion and then
simulates descent to impact under an end-of-flight family $h$ --- unpiloted glide after loss of
control, piloted glide, and a family in which control inputs continue --- producing the terminal
trajectory $\traj_{K:I}$ and, through $\Psi$, the impact complex $\impact$.

The two transmissions at 00:19 are the only evidence of this phase, and they are treated as an
enumerated choice rather than a decision. Six data treatments are defined, from holding the
00:19 data out entirely, through using the first burst timing offset alone, to using both
frequency offsets with and without the oscillator start-up correction of Holland (2018). The
anomalous second burst timing offset is used only in the reproduction configuration. Against
these are set two causes of the power interruption --- fuel exhaustion followed by auxiliary
power unit start, and an unattributed cause --- and the two are combined per trajectory as
competing explanations,
%
\begin{equation}
L_{19} \;=\; p\,f_{\mathrm{lag}} \;+\; \frac{1-p}{W},
\end{equation}
%
in which $f_{\mathrm{lag}}$ is the likelihood of the observed log-on time under an Erlang lag
following exhaustion, $W$ normalises a uniform alternative over the admissible window, and $p$
is the mixing weight on the exhaustion cause. Bayes factors are interpreted on the scale of Kass and Raftery (1995), and are quoted only
within a fixed data treatment, comparing causes, or within a fixed cause, comparing treatments: a factor computed
across different data sets would not be a comparison of models. Hard selection of trajectories
by whether they reach fuel exhaustion in the admissible window is not used, because it charges
nothing for an unexplained log-on.

\subsection*{3.6.\quad \textit{Stage 3: the impact-level evidence}}

Four modules act on the shared impact samples, each returning a likelihood of its own data
given the impact complex. To make them composable, each obeys five rules. It returns a
likelihood and not a posterior, having removed any internal reference prior. It consumes each
observation exactly once, and declares which. It evaluates on the common impact samples and does
not smooth privately over impact location, because the product of separately smoothed
likelihoods is not the smoothed product. It is smooth, with explicit model error and no
probability floors, and it either covers the support of the prior or declines to return a value
--- a value that was not computed is recorded as such, and is not treated as impossible. And it
is accurate where the other evidence has its mass rather than only near its own peak.

\emph{Ocean drift} propagates each impact position forward under an ocean reanalysis with
windage and Stokes drift marginalised, and scores the discovery place and time of each recovered
item. Two global ocean products are carried --- one of them the Copernicus global reanalysis of
Lellouche et al.\ (2021) --- and averaged exactly at equal prior rather than chosen,
%
\begin{equation}
L_{M} \;=\; \tfrac{1}{2}\bigl(L_{G} + L_{C}\bigr),
\end{equation}
%
with the posterior weight of a product reported rather than assumed:
%
\begin{equation}
P\bigl(m = G \mid \Ev\bigr) \;=\; \frac{Z_{G}}{Z_{G} + Z_{C}}.
\end{equation}

\emph{Unsuccessful search} contributes the classical disconfirming term. If $P_D(\impact)$ is
the probability that the searches would have detected the wreckage had it been at $\impact$,
then
%
\begin{equation}
p\bigl(\impact \mid \Ev^{S}, \Ev^{-}, \Ev^{19}\bigr)
 \;\propto\; \bigl[\,1 - P_D(\impact)\,\bigr]\;
   p\bigl(\impact \mid \Ev^{-}, \Ev^{19}\bigr),
\end{equation}
%
following the treatment of Stone et al.\ (2014) for Air France 447. $P_D$ is spatially varying
and depends on the extent and piece sizes of the wreckage field, not only on its centre, which
is why the sinking stage below carries samples rather than summary moments.

\emph{Satellite imagery} enters as a likelihood on the same samples and conditional on the
imaged objects being debris from this aircraft; it is reported as a conditional hypothesis, never
as a selection of trajectories that could reach it. \emph{Hydroacoustics} predicts the signal
level and detection probability at each monitoring station given the impact energy, its
deposition duration and the propagation path, and returns the likelihood of the observed record.

\subsection*{3.7.\quad \textit{Impact to seabed}}

Sinking is a transform and not a likelihood: there is essentially no observed seabed debris to
score, and the one scoreable fact --- that the searches found nothing --- belongs to the search
term of~(9). Each impact sample therefore fans out into wreckage samples whose weights sum to
the parent's,
%
\begin{equation}
\bigl(\impact_i,\ w_i\bigr) \;\longmapsto\;
  \bigl\{\bigl(\mathbf{r}_{i\ell},\ w_i\,\pi_{i\ell}\bigr)\bigr\}_{\ell=1}^{L},
  \qquad \sum_{\ell=1}^{L} \pi_{i\ell} = 1,
\end{equation}
%
where $\mathbf{r}_{i\ell}$ is the resting place of the $\ell$th element of the field. This
preserves the extent of the field, which is exactly what the detection probability in~(9)
depends on.

\subsection*{3.8.\quad \textit{Stage 4: composition, mixing and evidence}}

Composition multiplies the inherited weight of each impact sample by the enabled likelihood
factors:
%
\begin{equation}
w_i \;\propto\; w_i^{\mathrm{flight}+\mathrm{EoF}}\;
  L_{19}(\impact_i)\;\prod_{j \in \mathcal{J}} L_j(\impact_i),
  \qquad
  \mathcal{J} \subseteq \{D, S, P, A\},
\end{equation}
%
with the evidence of a composed product given by $Z_{\mathcal{J}} = \sum_i w_i$. Because
$\mathcal{J}$ is selectable, the contribution of each line of evidence is measured by the change
in $\ln Z$ and in the posterior when it is added, which is how the results of section~5 are
quoted.

Core strata are mixed, not pooled, with weights formed from an equal prior over families and the
evidence each family achieves:
%
\begin{equation}
W_s \;\propto\; P_{\mathrm{core}}(s)\, Z_s ,
\end{equation}
%
and results are reported per end-of-flight family first; any average across families appears
only beside the sensitivity to the family prior. This matters because a family prior is a
choice, and a result conditional on that choice must not be presented as data-driven.

\subsection*{3.9.\quad \textit{Convergence diagnostics}}

Two diagnostics are reported for every product. The effective sample size is computed per
evidence factor, so that a factor which concentrates the weight on a few rows is visible rather
than hidden in an aggregate. And replicate agreement is measured by split-half comparison: with
$R$ independent replicates, every balanced partition into two halves is formed, the posterior is
computed from each half, and the agreement between halves is quoted as a mean and a range over
all partitions. The threshold is replicate-count dependent, and at eight replicates a value of
%
\begin{equation}
\mathrm{SH} \;\geq\; 0.924
\end{equation}
%
is required before a distribution is called converged. Per-epoch effective sample size is not a
substitute: raising the worst epoch's value by a factor of four has been observed to change no
measure of replicate agreement. Where a product fails~(13) it is reported as not converged, and
the remedy is treated as a question for the sampler rather than for the smoothing.

% =====================================================================
\section*{4.\quad WHAT WAS CHANGED AND WHAT WAS ADDED}
% =====================================================================

\subsection*{4.1.\quad \textit{Fuel and endurance inside the likelihood}}

The published filter assumes infinite fuel and applies the fuel constraint afterwards as a
censor on trajectories, using broad fuel-burn information only to inform the allowable range of
Mach numbers (Davey et al., 2016, assumption~4). The fuel model here is an extension: burn is
integrated along each trajectory from the aircraft's known load, with two-tank bookkeeping and a
single engine-flameout event, and the resulting endurance enters the likelihood so that a
trajectory inconsistent with the fuel load is down-weighted while the filter advances. The
consequence is not merely bookkeeping. Where the published treatment produced a posterior and
then removed part of it, the integrated treatment changes which trajectories the proposal
explores, and therefore changes the shape of the posterior rather than only its support.

\subsection*{4.2.\quad \textit{Radar, waypoints and the sampler}}

Primary radar enters the likelihood rather than only the prior. A waypoint route family is added,
in which the aircraft is assumed to track between published navigation fixes, with turn
behaviour modelled as fly-by with leg tracking; this was adopted after a sensitivity test showed
that the alternative fly-over and direct-to model moved route log evidence by up to 7.6 nats and
the 00:19 latitude by about 0.7$^\circ$, which is too large to leave as an undeclared modelling
choice. The sampler is modified by tempering where the posterior separates into families. Four
of the five sampler interventions attempted did not improve convergence, and all four remain
runnable behind configuration flags: the pattern across them is itself a result, and is reported
in section~5.

\subsection*{4.3.\quad \textit{End of flight}}

The end-of-flight stage is the largest extension. It carries three descent families, a provisional
breakup model, and --- following evidence that piloted descents in an earlier configuration flew
implausibly straight --- a turn model in which descents may make random turns, a search sweep of
one to three circles, or a ditching approach, with bank angles of 10 to 30$^\circ$. Every choice
is recorded per descent so that the probability of a sweep, of a circuit, and of an into-wind as
against an along-swell final heading can be re-weighted afterwards from the evidence rather than
assumed in advance.

Two limits of the present implementation bound everything downstream of it and are stated here
rather than in a limitations paragraph. Commanded descent rates are capped, so the most rapid
descents consistent with the 00:19 frequency offsets are not reachable by the model; and free
flight is integrated at fixed lift coefficient without pitch dynamics, so uncontrolled push-over
states are not reachable either. In addition, the absence of a structural limit and of sufficient
transonic drag rise allows contact speeds near and above the speed of sound: between 11 and 21
per cent of the composed weight lies above 1.25 times maximum operating speed before the search
term is applied. Seabed and hydroacoustic quantities derived from those rows are therefore upper
bounds taken at the intact-impact point, not estimates.

\subsection*{4.4.\quad \textit{Drift of the recovered debris}}

This is an update rather than an extension. The published work modelled the drift of a single
recovered item; the treatment here uses the full set of recovered items, carries two global ocean
products averaged exactly at equal prior as in~(7), and marginalises over windage and Stokes
drift rather than fixing them. The authoritative drift work for this accident
(Griffin et al., 2016, 2017) is used for the leeway characterisation of the recovered flaperon
section, with the leeway and windage treatment following Sutherland et al.\ (2020) and the
debris-drift simulation strategy of Durgadoo et al.\ (2021). One term is specified and not yet built, and it is the most consequential omission in
the paper: the absence of recoveries on the Western Australian coast is not modelled, because the
beaching geometry covers only the coastlines on which items were found. Coasts outside those
boxes carry no identification, so the failure to recover debris in Western Australia cannot
presently down-weight impact locations whose debris would have beached there. Every result
involving drift carries that label.

\subsection*{4.5.\quad \textit{Unsuccessful search, imagery and hydroacoustics}}

The search term of~(9) is an extension: the published work discussed search areas without
estimating them. Detection probability is spatially varying and is presently evaluated at a
single conservative value of the detection parameter, which understates the strength of the
term; and the most recent search operation is excluded from the coverage used here. Both
qualifications work in the same direction, so the search term as reported is a lower bound on
its true effect.

The imagery term is a conditional hypothesis built on the official analysis of the Pl\'eiades 1A
imagery (Minchin et al., 2017). It is evaluated as a likelihood on the shared
impact samples conditional on the imaged objects being debris from this aircraft, and every
product involving it carries that condition. It is never implemented as a pre-selection of
trajectories that could reach the imaged position, because selecting trajectories by whether
they reach a hypothesised outcome is not inference.

The hydroacoustic term is predictive. It is built and it runs, but its transmission loss is not
absolutely calibrated, for five specific reasons set out in section~6, and its log-likelihood
ratios accordingly lie within noise.

\subsection*{4.6.\quad \textit{Verification and validation}}

An estimator assembled from independently developed parts needs its joins tested as much as its
physics, and four classes of test are used.

The reproduction is a regression test. A configuration holding the published priors and models
must reproduce the stored log evidence to the last digit, and the test is run after every change
to the filter. This is what separates the two classes of change that matter in practice: a change
to the filter's numerics --- a measurement model, a prior, the integration step --- invalidates
existing runs and obliges a re-run, whereas a change downstream of the filter leaves the
filter's output byte-identical and existing runs valid. Distinguishing the two has a direct cost
consequence, because a full-scale run is of the order of a day of wall-clock time on the hardware
used here.

Each module carries a synthetic-recovery test: data are generated from a known state, the module
is run, and the coverage of its posterior intervals is checked against nominal. This catches the
failure in which a likelihood is correct up to a constant but mis-scaled, which composition would
otherwise silently propagate. Each module also carries hand-computed fixtures for its core
quantities, so that a refactoring that changes a number is detected rather than absorbed.

The geodetic and satellite-geometry computations are verified against independent
implementations --- geodesic distances and azimuths against the algorithms of Karney (2013) and
the earlier method of Vincenty (1975) --- and the composition stage is verified by reproducing,
exactly, the quantities that the end-of-flight and sinking stages report independently for the
same rows. A composition that cannot reproduce its inputs is not combining evidence, it is
re-deriving it.

Finally, a vectorised re-computation of the 00:19 timing residual from stored final positions is
checked against a scalar reference on a random sample of rows at every run, with agreement
required to one part in $10^{6}$~$\mu$s. The re-computation matters beyond speed: because it is
performed from the saved positions rather than read from what the filter weighted, it provides an
out-of-sample check on the rungs of the evidence ladder that were not scored.

% =====================================================================
\section*{5.\quad RESULTS}
% =====================================================================

\subsection*{5.1.\quad \textit{What converged, and what did not}}

The reproduction of the published model is exact: a configuration holding the published priors
reproduces the stored log evidence with a difference of zero, and continues to do so after every
extension is added, which is what makes the comparisons below meaningful.

Nothing composed is converged. Against the eight-replicate split-half floor of 0.924 in~(13),
the four core families return 0.916 for the waypoint family, 0.891 for the reproduction family,
0.830 for free flight and 0.794 for the descent-and-climb family. Carried into the composed
products, the three balanced partitions of the replicates differ by 11 to 57 per cent in the
area of the 90 per cent region and by up to 0.98$^\circ$ in its median. The non-convergence is
established as a property of the flight-path run rather than of the composition, by the
per-replicate spread of log evidence within a stratum, which is 0.20 to 0.36 nats in the
waypoint family. The seabed distribution inherits it: the two halves of the replicate set differ
by 10 to 46 per cent in the area of the 90 per cent region.

This is reported first because it conditions everything that follows. The orderings and
directions below are stable across partitions; the areas and medians are not.

\subsection*{5.2.\quad \textit{Unsuccessful search is the largest single factor}}

Adding the search term removes 54 to 71 per cent of the probability mass, a change in log
evidence of $-0.78$ to $-1.23$. It moves the median impact latitude 0.3 to 1.2$^\circ$ north and
\emph{enlarges} the 90 per cent region by 60 to 96 per cent: under the treatment using the first
00:19 burst timing offset alone, from 165{,}000 to 309{,}000~km$^2$. The direction is the
informative part. A search that found nothing does not concentrate a distribution; it removes the
places that were searched and redistributes the mass to places that were not, which here lie
north of and outside the searched corridor. Reported under: not converged; detection parameter
evaluated at a single conservative value, understating the term; most recent search operation
excluded.

\subsection*{5.3.\quad \textit{Drift widens the distribution and does not discriminate}}

Drift of the recovered debris widens the flight and end-of-flight region rather than narrowing
it --- from 117{,}000 to 165{,}000~km$^2$ under the same 00:19 treatment --- and shifts it
north-east. Its evidence does not discriminate between the treatments of the 00:19 data, which
span a range in log evidence of 0.10 nats, from $-128.06$ to $-128.16$. Reported under: not
converged; Western Australia non-recovery not modelled; the drift factor is applied only to rows
that the term could be computed for.

Nor does the debris evidence select an ocean product. The posterior weight of the first product
in~(8) is 0.46 to 0.51 with drift alone, 0.47 to 0.51 after the search term, and 0.54 to 0.63
when the imagery term is included, so the absolute log Bayes factor never exceeds 0.53. On this
evidence the choice of global ocean reanalysis is not resolved by the debris, and averaging over
it is the correct treatment rather than a convenience.

\subsection*{5.4.\quad \textit{The 00:19 transmissions move the answer less than expected}}

Across the six data treatments, the median of the drift-composed product moves by at most
0.43$^\circ$, from $-36.86^\circ$ with the 00:19 data held out to $-37.29^\circ$ with both
frequency offsets and the oscillator start-up correction applied. Reported under: not converged;
conditional on an equal family prior; end-of-flight physics provisional, with a provisional
turn model, an interim pitch-down treatment and a provisional breakup model; Western Australia
non-recovery not modelled.

The ranking of treatments is not estimable for the cells that matter most. The two cells that
correspond to the named hypotheses in the literature fall below the thousand-effective-row floor
once the later evidence is applied: one falls from 275 effective rows to between 29 and 220 when
composed, the other from 833 to between 42 and 562, and every product using both frequency
offsets drops below the floor after the search term or under the imagery hypothesis. The
conclusion that follows is a negative one and is stated as such: with the present sampler, the
data treatment that the 00:19 evidence favours cannot be determined, and the apparent stability
of the median across treatments is partly a statement about the sampler rather than about the
evidence.

\subsection*{5.5.\quad \textit{Two negative results worth keeping}}

Adding a vertical rate to the cruise state --- an omission in the published model that appeared
likely to matter, because the 184~Hz step between the two 00:19 frequency offsets corresponds to
an equivalent vertical rate of about 10{,}500~ft/min, some nine times the largest bias excursion
in the published work --- changes the posterior median by 0.03$^\circ$ or less and the
probability of the secondary mode by about two percentage points. The omission is an order of
magnitude larger than the noise it was being scored against, and it is still not a lever on the
secondary mode, because the altitude changes it represents are too brief.

Four of the five sampler interventions attempted failed to improve replicate agreement, and the
one that succeeded did so by aiming the proposal at the cells where the families separate rather
than by adding particles. Raising the worst epoch's effective sample size by a factor of 4.3
changed no measure of convergence. The pattern is that this posterior's difficulty is a
structural multi-modality and not a shortage of samples.

\subsection*{5.6.\quad \textit{The seabed distribution}}

The seabed regions coincide with the impact regions to within 2.5 per cent for matched outcomes,
which is expected: the horizontal displacement during sinking is small against the width of the
impact distribution. Reported under: not converged; the seabed resample is re-weighted from the
composed weights rather than drawn from them, which is a partial treatment; and breakup rows are
settled from the intact impact point, so the extent of the field is a lower bound.
Figure~8 shows the seabed distribution beside the impact distribution it derives from.

\subsection*{5.7.\quad \textit{Evidence accounting, and where the samples go}}

Because the set $\mathcal{J}$ in~(11) is selectable, each line of evidence can be costed in two
currencies: the change it makes to the log evidence, and the effective sample size it leaves
behind. The two do not move together, and the difference is the most useful diagnostic in the
whole exercise.

The search term changes the log evidence by $-0.78$ to $-1.23$ nats --- a modest number --- while
removing most of the probability mass and reshaping the distribution substantially. The drift
term, by contrast, sits at around $-128$ nats, which reflects the improbability of the observed
discovery pattern under \emph{any} impact location rather than discrimination between locations:
the spread of that figure across the six 00:19 treatments is only 0.10 nats, and it is the spread,
not the level, that carries information. Reporting the level alone would invite the reading that
drift is the dominant evidence, when in fact it is the least discriminating term in the
composition. Figure~6 shows the evidence of each product with its replicate spread, which is the
form in which these numbers should be read.

Effective sample size tells the other half of the story. The treatment using the first burst
timing offset with the first frequency offset retains of the order of $6\times10^{4}$ effective
rows after flight and end of flight, $4\times10^{4}$ after drift, and $1\times10^{4}$ after the
search term, falling to a few thousand under the imagery hypothesis. That remains estimable. The
treatments using both frequency offsets do not: as reported in section~5.4, their named cells
fall to between a few tens and a few hundred effective rows once the later evidence is applied.
One line of evidence is therefore capable of making another unquotable, and the only honest
response is to report the effective sample size per factor beside every number, as in Figure~4.

A related accounting point is worth recording because it closes a gap that was open at the
previous pass. The rows for which the drift term could not be computed amount to at most 0.11 per
cent of the base weight in any stratum and treatment, and excluding them changes the log evidence
by less than 0.001 and moves no median or area visibly. At the previous composed pass, that
uncomputed fraction was large enough that whether it was carried or excluded changed the composed
areas materially. It no longer is, and the difference is a property of the run rather than of the
composition.

% =====================================================================
\section*{6.\quad DISCUSSION}
% =====================================================================

\subsection*{6.1.\quad \textit{What integration bought}}

Three of the results above could not have been obtained without carrying one distribution
through the whole chain. That unsuccessful search is the dominant factor is a statement about the
ratio of two terms in~(11), and is only available if both are evaluated on the same samples. That
drift widens rather than narrows the distribution contradicts the intuition that more evidence is
more constraint, and it has a clear cause: the drift likelihood is broad in the along-arc
direction and its peak lies away from the flight-path peak, so multiplying the two moves mass
outward. And the inability to rank the 00:19 treatments is visible only because the effective
sample size is tracked per factor rather than in aggregate.

The integration also exposes a dependency that separate analyses conceal. The detection
probability in~(9) depends on the extent and the piece sizes of the wreckage field, which depend
on the breakup state at impact, which depends on the descent family. A search term evaluated
against a point estimate of impact location therefore understates its own uncertainty. This is
the reason the sinking stage carries samples, as in~(10), rather than summary moments.

\subsection*{6.2.\quad \textit{The limits that bound the present answer}}

Four limits bound the result, in descending order of effect.

The flight-path run is not converged, and the remedy is a sampler aimed at the cells where the
families separate rather than more particles. Until it is converged, the areas and medians
reported here are not quotable to the precision at which they are printed.

The end-of-flight physics does not reach the whole feasible set. Commanded descent rates are
capped below the most rapid descents consistent with the 00:19 frequency offsets, and free
flight carries no pitch dynamics, so uncontrolled push-over states are unreachable. Both bound
the impact distribution's spread and both are reach limitations rather than coverage
limitations: the model cannot produce those states at all, so no amount of sampling would find
them.

The absence of recoveries in Western Australia is not modelled, which removes a line of
disconfirming evidence that would act on precisely the impact locations the search term pushes
mass towards.

And the absence of a structural limit allows impact speeds that no airframe would reach, so
quantities derived from the fastest rows are upper bounds.

\subsection*{6.3.\quad \textit{Hydroacoustics: what it is now, and what it could be}}

That a water impact of this kind is detectable in principle at long range has been established
by Metz et al.\ (2023), and the dependence of such a detection on propagation conditions and
impact parameters by Duncan and Dall'Osto (2023); the impact-generated acoustic and
acoustic-gravity wave physics is treated by Kadri (2024), and the record for this accident
specifically was examined by Duncan et al.\ (2014). Propagation here is modelled with the
standard methods (Jensen et al., 2011), using the split-step Pad\'e parabolic equation solution
of Collins (1993).

Even so, the hydroacoustic term as implemented is predictive and not evidential. Its log-likelihood ratios lie within
noise, and the reason is calibration rather than physics. Five specific gaps are responsible.
There is no near-surface calibration source of known spectrum, so the directivity of a seismic
air-gun source cannot be separated from its coupling into the water column. The deep explosive
source spectrum below the bubble frequency is unknown to between 4 and 17~dB, because the
standard source-level relations were established for shallower charges than the calibration shots
available. No frequency band below 16~Hz is calibrated by any deep source, and the bands in which
an impact marker would be expected lie below 12.5~Hz. At one of the two southern stations, the
Chagos Bank blocks every modelled path from every candidate impact location, so a non-detection
there carries no information at all. And at a bathymetric ridge cutting the sound-channel axis,
the modelled transmission loss is 49 to 55~dB too lossy, an error that two-dimensional
propagation modelling does not remove.

The quantitative consequence is worth stating precisely, because it is the difference between an
unused observation and a decisive one. An unarrested rapid descent deposits of the order of
several gigajoules at the surface, which is roughly an order of magnitude above the energy of the
nearest reference event for which a calibrated near-surface recording exists. If the
transmission loss were absolutely calibrated and the sub-16~Hz band characterised, the predicted
signal level at the unobstructed stations would separate the descent families by a margin that
the present uncertainty entirely masks --- and, equally informative, a confirmed non-detection
would down-weight the most energetic descents. The limit today is the power of the test, not the
model.

Five pieces of work would close those gaps, each tied to the gap it closes: access to the raw
hydrophone data rather than to derived products, through a computation-to-data arrangement;
absolute calibration from the geometry of a known source array together with its calibrated
near-source recordings; a near-surface source and coupling characterisation; three-dimensional
propagation across the ridges; and suppression of the seismic air-gun shot trains using survey
source and navigation logs. The first and second are the largest.

It should also be recorded that the bench-test evidence underlying the oscillator start-up
correction applied to the 00:19 frequency offsets is an unpublished internal study, available
only through the summary of it in Holland (2018). The correction is used here, and the
sensitivity to using it is reported, but a referee cannot inspect its primary evidence.

% =====================================================================
\section*{7.\quad CONCLUSIONS AND RECOMMENDATIONS}
% =====================================================================

An integrated Bayesian reconstruction of MH370 is feasible and has been built: one probability
distribution is carried from the last radar fix through the satellite handshakes, the end of
flight, water impact and sinking, and is re-weighted by ocean drift, unsuccessful search,
satellite imagery and hydroacoustic prediction on a single shared set of impact samples.

Within that reconstruction, unsuccessful seabed search is the largest single factor acting on the
impact distribution, removing 54 to 71 per cent of the probability and moving the remainder north
and outward. Drift of the recovered debris widens the distribution rather than narrowing it, and
does not discriminate between the available treatments of the final two transmissions, nor
between global ocean reanalyses. The final two transmissions move the median by less than half a
degree across all six treatments, but the cells corresponding to the named hypotheses in the
literature are not estimable with the present sampler, so that stability is in part a statement
about the sampler.

No composed distribution is converged against the replicate-agreement criterion proposed here.
The orderings and the directions of effect are stable; the areas and the medians are not. This is
stated as a result because the alternative --- reporting areas to three significant figures from
an unconverged run --- would be worse than reporting nothing.

Four lines of work follow, and they are ordered by how much they would change the answer rather
than by how difficult they are. First, a sampler aimed at the cells where the end-of-flight
families separate, which is the binding constraint on every number in this paper. Second, a
descent model that reaches the rapid descents the 00:19 frequency offsets admit, together with a
structural limit that excludes the impact speeds the present model allows. Third, the Western
Australian non-recovery term, which would act directly on the region the search term favours.
Fourth, absolute calibration of the hydroacoustic transmission loss and characterisation of the
sub-16~Hz band: of the four, this is the only one that could turn an unused observation into a
decisive one, and its value of information is correspondingly the largest.

% =====================================================================
\section*{ACKNOWLEDGEMENTS}
% =====================================================================

[TO BE COMPLETED BY PETE. Note that the journal must not be thanked, and that
acknowledgements precede the reference list.]

\section*{COMPETING INTERESTS}
Competing interests: The author declares none.

\section*{USE OF ARTIFICIAL INTELLIGENCE TOOLS}
[TO BE COMPLETED, AND REQUIRED. The journal requires a declaration describing any AI tool used
to generate text or images, or to analyse or extract insights from data, giving the tool name and
version, the dates of use, how others can access it, and a full description of the use. Analysis
use is declared in the methods section rather than in an acknowledgement. This estimator was
developed and run with AI assistance throughout, and the declaration must be assembled from what
the repository records rather than written from memory. The public repository, configurations and
commit history support a specific and verifiable statement. See paper/JON-RULES.md \S7.]

\section*{DATA AVAILABILITY}
[PROPOSED. The engine, configurations and result notes are in a public repository; the
aircraft-performance tables and the publisher's copy of the reference monograph are not
redistributable and are identified rather than included. Confirm against the journal's research
transparency policy, which has not yet been read.]

% =====================================================================
\section*{TABLES}
% =====================================================================

\begin{table}[htbp]
\centering
\caption{Elements of the integrated estimator, classified against the published reconstruction,
with build status. R = reproduction, U = update, M = modification, E = extension.
B = built and run, S = specified but not run.}
\vspace{4pt}
\small
\begin{tabular}{@{}p{4.9cm}cc p{6.6cm}@{}}
\toprule
Element & Class & Status & Note \\
\midrule
Satellite measurement models        & R & B & Ashton et al.\ (2015), unchanged \\
Cruise and manoeuvre dynamics       & R & B & Davey et al.\ (2016), ch.\ 6--7 \\
Baseline reproduction               & R & B & Reproduces stored log evidence exactly \\
Primary radar prior                 & R & B & Davey et al.\ (2016), ch.\ 4 \\
Ocean drift of recovered debris     & U & B & More items; two ocean products averaged \\
Composition across evidence         & U & B & Separate crate; not in the main runner \\
Prior track, sampler, 00:19 treatment & M & B & Tempering; six enumerated data treatments \\
Fuel and endurance in the likelihood & E & B & Two-tank; single flameout event \\
Radar in the likelihood             & E & B & Additional to the prior \\
Waypoint route family               & E & B & Fly-by with leg tracking \\
Vertical rate in the cruise BFO     & E & B & Negative result; see \S5.5 \\
End of flight to impact             & E & B & Three families; provisional breakup \\
Searched-area likelihood            & E & B & Lower bound; see \S4.5 \\
Impact to seabed (sinking)          & E & S & Transform; carried as samples \\
Hydroacoustics                      & E & S & Predictive only; see \S6.3 \\
Satellite imagery hypothesis        & E & B & Conditional on objects being debris \\
Western Australia non-recovery      & E & S & Not modelled; see \S4.4 \\
\bottomrule
\end{tabular}
\end{table}

\begin{table}[htbp]
\centering
\caption{Treatments of the 00:19 transmissions, and the causes against which they are set. Bayes
factors are quoted only within a row or within a column, never across both.}
\vspace{4pt}
\small
\begin{tabular}{@{}p{1.6cm}p{8.6cm}p{4.4cm}@{}}
\toprule
Treatment & Data used & Cause \\
\midrule
D0 & 00:19 held out entirely & C1: fuel exhaustion followed \\
D1 & First burst timing offset only & \quad by auxiliary power unit \\
D2 & D1 + first frequency offset as recorded & \quad start, with an Erlang lag \\
D3 & D1 + both frequency offsets as recorded & \\
D4 & D3 with the start-up correction on both & C2: unattributed, uniform \\
D5 & D3 with the correction on the second only & \quad over the admissible window \\
\bottomrule
\end{tabular}
\end{table}

\begin{table}[htbp]
\centering
\caption{Coverage limitations that bound the reported distributions, with the direction of their
effect. A reach limitation means the model cannot produce the states concerned, so no amount of
sampling would find them.}
\vspace{4pt}
\small
\begin{tabular}{@{}p{5.4cm}p{2.2cm}p{7.0cm}@{}}
\toprule
Limitation & Kind & Effect on the reported result \\
\midrule
Flight-path run not converged & coverage & Areas and medians not quotable at printed precision \\
Descent rates capped & reach & Rapid descents admitted by 00:19 unreachable \\
No pitch dynamics in free flight & reach & Uncontrolled push-over states unreachable \\
No structural limit & reach & Impact speeds above 1.25 $V_{MO}$ on 11--21\% of weight \\
Western Australia non-recovery & reach & Disconfirming evidence absent where mass is moved to \\
Detection parameter conservative & coverage & Search term is a lower bound \\
Transmission loss uncalibrated & coverage & Hydroacoustic ratios within noise \\
\bottomrule
\end{tabular}
\end{table}

% =====================================================================
\section*{REFERENCES}
% =====================================================================
\begingroup
\setlength{\parskip}{3pt}
\small

Ashton, C., Shuster Bruce, A., Colledge, G. and Dickinson, M. (2015). The Search for MH370.
\textit{The Journal of Navigation}, 68, 1--22. https://doi.org/10.1017/S037346331400068X.

Australian Transport Safety Bureau. (2017). The Operational Search for MH370. ATSB,
AE-2014-054.

Collins, M. D. (1993). A split-step Pad\'e solution for the parabolic equation method.
\textit{Journal of the Acoustical Society of America}, 93, 1736--1742.
https://doi.org/10.1121/1.406739.

Davey, S., Gordon, N., Holland, I., Rutten, M. and Williams, J. (2016). \textit{Bayesian
Methods in the Search for MH370}. Springer. https://doi.org/10.1007/978-981-10-0379-0.

Duncan, A. J. and Dall'Osto, D. R. (2023). Long-range underwater acoustic detection of aircraft
surface impacts --- the influence of acoustic propagation conditions and impact parameters.
\textit{Journal of the Acoustical Society of America}, 154, A49.
https://doi.org/10.1121/10.0022761.

Duncan, A. J., Gavrilov, A. N. and McCauley, R. D. (2014). Analysis of low frequency underwater
acoustic signals possibly related to the loss of MH370. Centre for Marine Science and
Technology, Curtin University.

Durgadoo, J. V., Biastoch, A., New, A. L., R\"uhs, S., Nurser, A. J. G., Drillet, Y. and
Bidlot, J.-R. (2021). Strategies for simulating the drift of marine debris. \textit{Journal of
Operational Oceanography}, 14, 1--12. https://doi.org/10.1080/1755876X.2019.1602102.

Griffin, D. A., Oke, P. R. and Jones, E. M. (2016). The search for MH370 and ocean surface
drift. CSIRO Oceans and Atmosphere, EP167888. https://doi.org/10.4225/08/5892224dec08c.

Griffin, D. A., Oke, P. R. and Jones, E. M. (2017). The search for MH370 and ocean surface
drift --- Part II. CSIRO Oceans and Atmosphere, EP172633.

Holland, I. D. (2018). MH370 Burst Frequency Offset Analysis and Implications on Descent Rate
at End-of-Flight. \textit{IEEE Aerospace and Electronic Systems Magazine}, 33, 24--33.
https://doi.org/10.1109/MAES.2018.170048.

Jensen, F. B., Kuperman, W. A., Porter, M. B. and Schmidt, H. (2011). \textit{Computational
Ocean Acoustics}. Springer. https://doi.org/10.1007/978-1-4419-8678-8.

Kadri, U. (2024). Underwater acoustic analysis reveals unique pressure signals associated with
aircraft crashes in the sea: revisiting MH370. \textit{Scientific Reports}, 14, 10102.
https://doi.org/10.1038/s41598-024-60529-1.

Karney, C. F. F. (2013). Algorithms for geodesics. \textit{Journal of Geodesy}, 87, 43--55.
https://doi.org/10.1007/s00190-012-0578-z.

Kass, R. E. and Raftery, A. E. (1995). Bayes Factors. \textit{Journal of the American
Statistical Association}, 90, 773--795. https://doi.org/10.1080/01621459.1995.10476572.

Lellouche, J.-M. and others (2021). The Copernicus Global 1/12$^\circ$ Oceanic and Sea Ice
GLORYS12 Reanalysis. \textit{Frontiers in Earth Science}, 9, 698876.
https://doi.org/10.3389/feart.2021.698876.

Malaysian ICAO Annex 13 Safety Investigation Team for MH370. (2018). Safety Investigation
Report: MH370 (9M-MRO). Ministry of Transport Malaysia, MH370/01/2018.

Metz, D., Obana, K. and Fukao, Y. (2023). Remote hydroacoustic detection of an airplane crash.
\textit{Pure and Applied Geophysics}, 180, 1343--1351.
https://doi.org/10.1007/s00024-022-03117-6.

Minchin, S., Mueller, N., Lewis, A., Byrne, G. and Tran, M. (2017). Summary of imagery analyses
for non-natural objects in support of the search for Flight MH370: Results from the analysis of
imagery from the PLEIADES 1A satellite undertaken by Geoscience Australia. Geoscience
Australia, Record 2017/013. https://doi.org/10.11636/record.2017.013.

Ng, H.-F., Zhong, Y., Groves, P. D. and Hsu, L.-T. (2024). Grid-based 3DMA GNSS with clustering
and Doppler velocity. \textit{The Journal of Navigation}, 77, 658--676. [VERIFY DOI]

Stone, L. D., Keller, C. M., Kratzke, T. M. and Strumpfer, J. P. (2014). Search for the
Wreckage of Air France Flight AF 447. \textit{Statistical Science}, 29, 69--80.
https://doi.org/10.1214/13-STS420.

Sutherland, G., Soontiens, N., Dufois, F. and others (2020). Evaluating the leeway coefficient
of ocean drifters using operational marine environmental prediction systems. \textit{Journal of
Atmospheric and Oceanic Technology}. https://doi.org/10.1175/JTECH-D-20-0013.1.
[AUTHOR LIST TO BE COMPLETED]

Vincenty, T. (1975). Direct and inverse solutions of geodesics on the ellipsoid with application
of nested equations. \textit{Survey Review}, 23, 88--93.
https://doi.org/10.1179/sre.1975.23.176.88.

\endgroup

\end{document}
